\chapter{Stroke Imaging}

\section*{Patient focus}
This chapter explains the examination or clinical problem in accessible language. Patients should use it to prepare questions and understand the purpose, experience, benefits, and limitations of imaging; local instructions and the official report take priority.


\label{ch:stroke_imaging}

Stroke is a time-sensitive medical emergency where "time is brain." Acute stroke is classified as either \textbf{ischemic} (85\%) due to thromboembolic occlusion, or \textbf{hemorrhagic} (15\%) due to vascular rupture. The primary goal of emergency stroke imaging is to exclude intracranial hemorrhage, identify salvageable brain tissue (penumbra), and locate large vessel occlusions (LVO) for intervention.

---

\section*{1. The Hyperacute CT Stroke Protocol}
The standard emergency workup for suspected acute stroke is a rapid multi-modal CT protocol:

\subsection*{Non-Contrast Head CT (NCCT)}
The first-line scan, completed in under 1 minute.
\textbullet \textbf{Exclude Hemorrhage}: Acute hemorrhage appears bright white (hyperdense). Excluding hemorrhage is the absolute first step, as it is a strict contraindication to intravenous thrombolysis (tPA).
\textbullet \textbf{Identify Early Ischemic Changes}:
  \textbullet *Hyperdense MCA Sign*: A bright white line in the Sylvian fissure representing an acute thrombus in the middle cerebral artery.
  \textbullet *Loss of Grey-White Differentiation*: Blurring of the lentiform nucleus margins or loss of the "insular ribbon sign" within hours of occlusion.
\textbullet \textbf{ASPECTS (Alberta Stroke Program Early CT Score)}: A 10-point quantitative score assessing early ischemic changes in the MCA territory. Every affected region (out of 10 defined areas) subtracts 1 point:
  \textbullet ASPECTS of \textbf{10} is normal.
  \textbullet ASPECTS \textbf{$< 6$} indicates a large established infarct core and a higher risk of hemorrhagic conversion if treated with thrombolysis.

\subsection*{CT Angiography (CTA) of the Head and Neck}
Performed with IV contrast from the aortic arch to the skull vertex.
\textbullet \textbf{LVO Detection}: Identifies cut-offs in major arteries (terminal internal carotid, MCA M1/M2 segments, or basilar artery) that make the patient eligible for mechanical endovascular thrombectomy (EVT).
\textbullet \textbf{Collaterals}: Evaluates retrograde blood flow to the ischemic territory; good collaterals preserve the penumbra longer.

\subsection*{CT Perfusion (CTP)}
Calculates tissue blood flow parameters using continuous scanning during contrast wash-in:
\textbullet \textbf{Infarct Core} (Dead Tissue): Defined as areas with severely reduced Cerebral Blood Flow (CBF $< 30\%$).
\textbullet \textbf{Ischemic Penumbra} (Salvageable Tissue): Defined as areas with delayed blood flow (Time-to-maximum tracer arrival, Tmax $> 6$ seconds) but preserved Cerebral Blood Volume (CBV).
\textbullet \textbf{Mismatch Ratio}: A large penumbra surrounding a small core (mismatch) indicates that revascularization will rescue brain tissue, extending the thrombectomy window up to 24 hours.

---

\section*{2. The Hyperacute MRI Stroke Protocol}
In patients with unknown stroke onset time (e.g. "wake-up stroke"), MRI is the preferred modality using the \textbf{DWI-FLAIR Mismatch}:
\textbullet \textbf{DWI (Diffusion-Weighted Imaging)}: Detects cytotoxic edema within minutes of stroke onset, appearing bright white.
\textbullet \textbf{FLAIR (Fluid-Attenuated Inversion Recovery)}: Takes 3--4 hours for ischemic changes to become bright.
\textbullet \textbf{DWI-FLAIR Mismatch}: A lesion that is \textbf{positive on DWI but negative on FLAIR} indicates hyperacute ischemia ($< 4.5$ hours old). This confirms the stroke occurred recently enough for the patient to safely receive IV thrombolysis despite the unknown onset time.

---

\section*{3. Structured Reporting Example: Acute Ischemic Stroke with LVO}

\begin{testbox}[Indications:]
74-year-old female presenting with acute right-sided hemiplegia and aphasia. Last known well 2 hours ago.
\end{testbox}
\begin{refbox}[Findings:]
\textbf{NCCT Head}: No intracranial hemorrhage or mass effect. There is subtle loss of grey-white differentiation in the left lentiform nucleus and insular cortex. The ASPECTS is calculated as 8 (loss of L-insular ribbon and L-lentiform nucleus). No hyperdense MCA sign.\\
\textbf{CTA Head \& Neck}: There is an abrupt cut-off of contrast enhancement in the M1 segment of the left middle cerebral artery, consistent with a large vessel occlusion. The bilateral internal carotid arteries and vertebral arteries are patent. Moderate retrograde collateral flow is visualized distal to the occlusion.\\
\textbf{CT Perfusion (Automated Software Analysis)}:
\textbullet Estimated Infarct Core (CBF $< 30\%$): 15 mL.
\textbullet Estimated Penumbra (Tmax $> 6$s): 85 mL.
\textbullet Mismatch Volume: 70 mL (Mismatch Ratio: 5.6).
\end{refbox}
\begin{pearlbox}[Impression:]
1. Acute occlusion of the left MCA M1 segment.
2. Moderate early ischemic changes in the left MCA territory (ASPECTS 8).
3. Large core-penumbra mismatch (70 mL salvageable tissue), highly favorable for endovascular mechanical thrombectomy.
\end{pearlbox}

\clearpage
